% !TeX program = xelatex
% Plantilla Pandoc del manual de uso de Mástil Escalas.
% Sigue la composición de docs/latex/DTS_Mastil_Escalas.tex: clase report, mismas
% fuentes, márgenes, títulos numerados, cabeceras, pies e índice.
\documentclass[a4paper,11pt,oneside]{report}

% ---------------------------------------------------------------------------
% Fuentes y unicode
% ---------------------------------------------------------------------------
\usepackage{fontspec}
\IfFontExistsTF{TeX Gyre Pagella}{\setmainfont{TeX Gyre Pagella}}{\setmainfont{P052}}
\IfFontExistsTF{TeX Gyre Heros}{\setsansfont{TeX Gyre Heros}}{\setsansfont{Nimbus Sans}}
\IfFontExistsTF{TeX Gyre Cursor}{\setmonofont{TeX Gyre Cursor}[Scale=0.92]}{\setmonofont{Nimbus Mono PS}[Scale=0.92]}

\usepackage[spanish,es-noshorthands]{babel}
\usepackage{microtype}

% ---------------------------------------------------------------------------
% Geometría de página
% ---------------------------------------------------------------------------
\usepackage[a4paper,margin=2.6cm,headheight=23pt,headsep=10pt,footskip=1.1cm]{geometry}
\raggedbottom
\setlength{\emergencystretch}{3em}
\hyphenpenalty=10000
\exhyphenpenalty=10000
\clubpenalty=10000
\widowpenalty=10000
\displaywidowpenalty=10000
\brokenpenalty=10000

% ---------------------------------------------------------------------------
% Color
% ---------------------------------------------------------------------------
\usepackage{xcolor}
\definecolor{navybg}{HTML}{0B1D2E}
\definecolor{graphitebg}{HTML}{1B2027}
\definecolor{navydeep}{HTML}{123049}
\definecolor{navytext}{HTML}{123049}
\definecolor{navylink}{HTML}{144870}
\definecolor{lightblue}{HTML}{8FD3FE}
\definecolor{amber}{HTML}{F5A623}
\definecolor{amberlink}{HTML}{7A4A00}
\definecolor{softgray}{HTML}{5B6570}

% ---------------------------------------------------------------------------
% Estructura, títulos
% ---------------------------------------------------------------------------
\usepackage{titlesec,needspace}
\setcounter{secnumdepth}{3}
\setcounter{tocdepth}{2}

\titleformat{\chapter}[display]
  {\normalfont\sffamily\huge\bfseries\color{navytext}}
  {\makebox[\textwidth][r]{\small\sffamily\color{softgray}CAPÍTULO \thechapter}}
  {6pt}
  {\Huge\raggedright}
\titlespacing*{\chapter}{0pt}{0pt}{28pt}

\titleformat{\section}
  {\normalfont\sffamily\Large\bfseries\color{navytext}\raggedright}
  {\thesection}{10pt}{}
\titlespacing*{\section}{0pt}{18pt}{8pt}

\titleformat{\subsection}
  {\normalfont\sffamily\large\bfseries\color{navytext}\raggedright}
  {\thesubsection}{8pt}{}
\titlespacing*{\subsection}{0pt}{14pt}{6pt}

% Índice: los capítulos llevan también puntos guía hasta el número de página.
% Es la definición de report.cls con los puntos de \@dottedtocline; se usa \(…\)
% porque el signo de dólar delimita las variables de la plantilla Pandoc.
\makeatletter
\renewcommand*\l@chapter[2]{%
  \ifnum \c@tocdepth >\m@ne
    \addpenalty{-\@highpenalty}%
    \vskip 1.0em \@plus\p@
    \setlength\@tempdima{1.5em}%
    \begingroup
      \parindent \z@ \rightskip \@pnumwidth
      \parfillskip -\@pnumwidth
      \leavevmode \bfseries
      \advance\leftskip\@tempdima
      \hskip -\leftskip
      #1\nobreak
      {\normalfont\leaders\hbox{\(\m@th\mkern \@dotsep mu\hbox{.}\mkern \@dotsep mu\)}\hfill}%
      \nobreak\hb@xt@\@pnumwidth{\hss #2}\par
      \penalty\@highpenalty
    \endgroup
  \fi}
\makeatother

% ---------------------------------------------------------------------------
% Cabeceras y pies
% ---------------------------------------------------------------------------
\usepackage{fancyhdr}
\pagestyle{fancy}
\fancyhf{}
\fancyhead[C]{\small\sffamily\color{softgray}Mástil Escalas --- $headtitle$}
\fancyfoot[L]{\small\sffamily\color{softgray}Aplicación $appversion$ · $edition$}
\fancyfoot[R]{\small\sffamily\color{softgray}\thepage}
\renewcommand{\headrulewidth}{0.4pt}
\renewcommand{\headrule}{\hbox to\headwidth{\color{navydeep!40!white}\leaders\hrule height \headrulewidth\hfill}}
\renewcommand{\footrulewidth}{0pt}
\fancypagestyle{plain}{
  \fancyhf{}
  \fancyfoot[L]{\small\sffamily\color{softgray}Aplicación $appversion$ · $edition$}
  \fancyfoot[R]{\small\sffamily\color{softgray}\thepage}
  \renewcommand{\headrulewidth}{0pt}
}

% ---------------------------------------------------------------------------
% Tablas, listas
% ---------------------------------------------------------------------------
\usepackage{array,longtable,booktabs,calc}
\newcounter{none}
\renewcommand{\arraystretch}{1.12}
\usepackage{enumitem}
\setlist[itemize]{leftmargin=1.4em,itemsep=2pt,topsep=4pt}
\setlist[enumerate]{leftmargin=1.6em,itemsep=2pt,topsep=4pt}
\providecommand{\tightlist}{\setlength{\itemsep}{2pt}\setlength{\parskip}{0pt}}

% ---------------------------------------------------------------------------
% Código, rutas, texto monoespaciado
% ---------------------------------------------------------------------------
\usepackage{fvextra,framed,seqsplit}
\fvset{fontsize=\small,breaklines=true,breakanywhere=true,frame=single,rulecolor=\color{navydeep!35!white},framesep=6pt}
\RecustomVerbatimEnvironment{verbatim}{Verbatim}{breaklines,breakanywhere,samepage=true,fontsize=\small}
\let\originaltexttt\texttt
\renewcommand{\texttt}[1]{\originaltexttt{\seqsplit{#1}}}

% ---------------------------------------------------------------------------
% Imágenes, diagramas, símbolos
% ---------------------------------------------------------------------------
% Pandoc 3.2+ envuelve las imágenes en \pandocbounded; con versiones anteriores no existe y esta definición no cambia nada.
\providecommand{\pandocbounded}[1]{#1}
\usepackage{graphicx,xurl,tikz}
\usetikzlibrary{calc}
\usepackage[section]{placeins}
\usepackage{amssymb,newunicodechar}

% Capturas: figuras flotantes que genera scripts/generateManual.mjs. Cada captura lleva un
% marco fino (la interfaz tiene fondos claros que, sin él, se confunden con la página) y un
% pie numerado a mano en el Markdown, por eso se usa \caption* sin etiqueta automática.
\usepackage{caption}
\captionsetup{font={small,sf},labelformat=empty,justification=centering,skip=6pt}
\setlength{\intextsep}{12pt}
\setlength{\textfloatsep}{14pt}
\newcommand{\manualcapture}[2]{%
  {\setlength{\fboxsep}{0pt}\setlength{\fboxrule}{0.4pt}%
   \fcolorbox{navydeep!25}{white}{\includegraphics[width=\dimexpr#1-0.8pt\relax]{#2}}}}

% Recuadros «Consejo» y «Atención»: en el Markdown son citas (> **Consejo.** …).
\renewenvironment{quote}{%
  \def\FrameCommand{{\color{amber}\vrule width 3pt}\fboxsep=7pt\colorbox{amber!9}}%
  \MakeFramed{\advance\hsize-\width\FrameRestore}\small\noindent\ignorespaces}%
  {\endMakeFramed}

% ---------------------------------------------------------------------------
% Enlaces
% ---------------------------------------------------------------------------
\usepackage[colorlinks=true,linkcolor=navylink,citecolor=navylink,urlcolor=amberlink,
  bookmarksnumbered=true,bookmarksopen=true,breaklinks=true,
  pdftitle={Mástil Escalas - $headtitle$},
  pdfauthor={Jesús Quevedo Rodríguez},
  pdfsubject={$headtitle$}]{hyperref}
\urlstyle{same}
\usepackage{bookmark}
\newunicodechar{♭}{\texorpdfstring{\ensuremath{\flat}}{\unichar{9837}}}
\newunicodechar{♯}{\texorpdfstring{\ensuremath{\sharp}}{\unichar{9839}}}
\newunicodechar{✓}{\texorpdfstring{\ensuremath{\checkmark}}{\unichar{10003}}}
\newunicodechar{✗}{\texorpdfstring{\ensuremath{\times}}{\unichar{10007}}}
\newunicodechar{→}{\texorpdfstring{\ensuremath{\rightarrow}}{\unichar{8594}}}
\newunicodechar{←}{\texorpdfstring{\ensuremath{\leftarrow}}{\unichar{8592}}}
\newunicodechar{Δ}{\texorpdfstring{\ensuremath{\Delta}}{\unichar{916}}}
\renewcommand{\topfraction}{0.95}
\renewcommand{\bottomfraction}{0.95}
\renewcommand{\textfraction}{0.05}
\renewcommand{\floatpagefraction}{0.6}
\setcounter{topnumber}{3}
\setcounter{totalnumber}{4}

% ---------------------------------------------------------------------------
\begin{document}

% =============================================================================
% PORTADA
% =============================================================================
\begin{titlepage}
\newgeometry{margin=0cm}
\thispagestyle{empty}
\begin{tikzpicture}[remember picture,overlay]
  % Fondo
  \fill[navybg] (current page.south west) rectangle (current page.north east);
  \shade[top color=navybg,bottom color=graphitebg]
    ([yshift=9.5cm]current page.south west) rectangle (current page.south east);

  % --- Mástil estilizado (banda inferior): cuerdas, trastes e incrustaciones ---
  % En la plantilla Pandoc «$$» produce un signo de dólar literal (sintaxis calc de TikZ).
  \begin{scope}
    \clip (current page.south west) rectangle ([yshift=10cm]current page.south east);
    \begin{scope}[shift={($$(current page.south west)+(-0.4,4.3)$$)},rotate=-5,scale=0.75]
      \draw[lightblue,opacity=0.35,line width=1pt] (0,0) rectangle (30,4.2);
      \foreach \x in {2.4215,4.7057,6.8637,8.9008,10.822,12.633,14.347,15.965,17.489,18.928,20.287,21.570,22.771,23.908,25.000,26.024,26.984,27.885,28.746}{
        \draw[white,opacity=0.22,line width=0.5pt] (\x,0) -- (\x,4.2);
      }
      \foreach \y in {0.35,1.05,1.75,2.45,3.15,3.85}{
        \draw[white,opacity=0.30,line width=0.35pt] (0,\y) -- (30,\y);
      }
      \foreach \x in {5.785,9.861,13.490,16.727,24.454,26.504}{
        \filldraw[lightblue,opacity=0.85] (\x,2.1) circle (0.13);
      }
      \filldraw[lightblue,opacity=0.85] (20.929,1.4) circle (0.12);
      \filldraw[lightblue,opacity=0.85] (20.929,2.8) circle (0.12);
    \end{scope}
  \end{scope}

  % --- Antetítulo -----------------------------------------------------------
  \node[anchor=north] at ([yshift=-3.1cm]current page.north) {%
    \sffamily\bfseries\fontsize{12}{14}\selectfont\color{amber}%
    \MakeUppercase{Proyecto Mástil Escalas}%
  };

  % --- Título ---------------------------------------------------------------
  \node[anchor=north,text width=17cm,align=center] at ([yshift=-4.1cm]current page.north) {%
    \sffamily\bfseries\fontsize{29}{35}\selectfont\color{white}%
    \hyphenpenalty=10000\exhyphenpenalty=10000\uchyph=0
    $covertitle$%
  };

  % --- Subtítulo ------------------------------------------------------------
  \node[anchor=north,text width=14.5cm,align=center] at ([yshift=-6.9cm]current page.north) {%
    \sffamily\fontsize{12.5}{16}\selectfont\color{lightblue!85!white}%
    \hyphenpenalty=10000\exhyphenpenalty=10000\uchyph=0
    $coversubtitle$%
  };

  % --- Ficha inferior -------------------------------------------------------
  \node[anchor=south west,text width=13.5cm,align=left] at ([xshift=1.9cm,yshift=0.9cm]current page.south west) {%
    \sffamily\fontsize{10.5}{15}\selectfont\color{white!92!black}
    \textcolor{amber}{\bfseries Autor:} Jesús Quevedo Rodríguez\\[2pt]
    \textcolor{amber}{\bfseries Aplicación:} versión $appversion$\\[2pt]
    \textcolor{amber}{\bfseries Fecha:} $edition$\\[2pt]
    \textcolor{amber}{\bfseries Repositorio:} \url{https://github.com/a01653/mastil_escalas}%
  };
\end{tikzpicture}
\end{titlepage}
\restoregeometry
\clearpage

% =============================================================================
% ÍNDICE
% =============================================================================
\pagenumbering{roman}
\tableofcontents
\clearpage
\pagenumbering{arabic}
\setcounter{page}{1}

$body$
\end{document}
